\documentclass[10pt]{article}

\usepackage[utf8]{inputenc}

\usepackage[T1]{fontenc}

\usepackage{amsmath}

\usepackage{amsfonts}

\usepackage{amssymb}

\usepackage[version=4]{mhchem}

\usepackage{stmaryrd}

\usepackage{mathrsfs}



\begin{document}

компоненты поля; принята система отсчета, в к-рой $\hbar=c=1)$. Тогда Н.т. утверждает, что из инвариантности действия (2) относительно преобразований с параметрами $\varepsilon^{A}$



\[

\begin{gathered}

x^{\mu} \rightarrow x^{\prime \mu}=x^{\mu}+\Lambda_{A}^{\mu}(x) \varepsilon^{A}, \\

\varphi^{a}(x) \rightarrow \varphi^{\prime a}\left(x^{\prime}\right)=\varphi^{a}(x)+\mathscr{L}_{A}^{a b}(x) \varphi_{b}(x) \varepsilon^{A}

\end{gathered}

\]



для произвольной области интегрирования $R$ вытекает дифференциальный закон сохранения:



\[

d J_{A}^{\mu} / d x^{\mu}=0,

\]



где так наз. н е т е р о в т ок $J_{A}^{\mu}$ вычисляется из лагранжиана по правилу:



\[

J_{A}^{\mu}=T_{v}^{\mu} \Lambda_{A}^{v}+\frac{\partial L(x)}{\partial \varphi_{, \mu}^{a}} \varphi^{b} \mathscr{L}_{A}^{a b},

\]



где



\[

T_{v}^{\mu}=\delta_{v}^{\mu} L(x)-\frac{\partial L(x)}{\partial \varphi^{a}} \varphi_{v}^{a}

\]



( $\delta_{v}^{\mu}$ - символ Кронекера; по повторяюшемуся индексу предполагается суммирование). Интегрируя (3) по произвольному 4-объему и используя теорему Гаусса, получаем, что полный 4-поток вектора $J_{A}^{\mu}$ через ограничивающую этот объем гиперповерхность равен нулю. Выбирая гиперповерхность в виде цилиндра с пространственноподобными основаниями, такого, что потоком через боковые стенки можно пренебречь, приходим к утверждению, что направленные в будущее потоки вектора $J_{A}^{\mu}$ через нижнее и верхнее основания равны. Отсюда следует, что нетеровы заряды



\[

Q_{A}(t)=\int_{x^{0}=t} J_{A}^{0}(x) d^{3} x,

\]



во-первых, сохраняются во времени (интегральная фо р м а Н.т.), во-вторых, преобразуются при Лоренца преобразованиях контравариантно соответствующим параметрам $\varepsilon^{A}$.



Из физич. представлений об однородности и изотропии пространства-времени следует, что для любой замкнутой системы действие должно быть инвариантно относительно преобразований Пуанкаре группы, что в силу Н. т. приводит к существованию 10 фундаментальных сохраняющихся величин: энергии, 3 компонент импульса и 6 компонент 4-момента. Сохранение энергии и импульса следует из инвариантности относительно трансляций $\delta x^{\mu}=a^{\mu}$. При этом $A=\mu, \mathscr{L}_{A}^{a b}=0$, нетеровы токи исчерпываются выражением (5) и образуют тензор энерги и-импульса. Сохраняющиеся «заряды» суть компоненты 4-импульса:



\[

P_{\mu}=\int_{x^{0}=t} T_{\mu}^{0} d^{3} x

\]



Из инвариантности относительно трех пространственных поворотов и трех преобразований Лоренца



\[

A=(\mu, v) ; \varepsilon^{A}=\omega^{v \mu}=-\omega^{\mu v} ; \Lambda_{A}^{\rho}=\delta_{v}^{\rho} g_{\mu \sigma} x^{\sigma}

\]



(где $g_{\mu \sigma}-$ метрич. тензор) вытекает дифференциальный закон сохранения для тензора плотности момента:



\[

M_{v \rho}^{\mu}=-\frac{1}{2} T_{v}^{\mu} x_{\rho}+\frac{1}{2} T_{\rho}^{\mu} x_{v}+\frac{\partial L(x)}{\partial \varphi_{, \mu}^{a}} \mathscr{L}_{v \rho}^{a b} \varphi^{b}

\]



$\mathscr{L}_{v \rho}^{a b}$ определяется спином полей. Соответствующий нетеров заряд есть 4-момент.



В гамильтоновом описании 10 фундаментальных величин являются генераторами соответствующих преобразований группы Пуанкаре и образуют относительно скобок Пуассона замкнутую алгебру Ли



\[

\begin{gathered}

\left(P_{\mu}, P_{v}\right)=0 ; \\

\left(M_{\mu \rho}, P_{v}\right)=-\left(g_{\mu v} g_{\rho \sigma}-g_{\mu \sigma} g_{\rho \mu}\right) P^{\sigma} ;

\end{gathered}

\]



\[

\left(M_{\mu \rho}, M_{v \sigma}\right)=-g_{\mu v} M_{\rho \sigma}+g_{\rho v} M_{\mu \sigma}+g_{\mu \sigma} M_{\rho v}-g_{\rho \sigma} M_{\mu v}

\]



изоморфную алгебре Ли группы Пуанкаре. Требование выполнения соотношений (7) в гамильтоновом формализме эквивалентно требованию инвариантности лагранжиана относительно группы Пуанкаре в лагранжевом формализме.



При наличии в системе симметрий, не связанных с пространством-временем, Н.т. позволяет построить и другие сохраняющиеся величины. При этом в выражении (4) для нетерова тока остается только второй член. Напр., если в системе с комплексным полем $\varphi^{a}$ действие инвариантно относительно глобального (с фазой $\alpha$, не зависящей от $x$ ) калибровочного преобразования 1-го рода



\[

\varphi^{a} \rightarrow \varphi^{a} \exp (i \alpha) ; \varphi^{* a} \rightarrow \varphi^{* a} \exp (-i \alpha),

\]



то будут сохраняться ток



\[

j^{\mu}(x)=i\left(\frac{\partial L(x)}{\partial \varphi_{, \mu}^{a}} \varphi^{a}-\varphi^{* a} \frac{\partial L(x)}{\partial \varphi_{, \mu}^{* a}}\right)

\]



и соответствующий заряд. В построении современных реалистич. квантовополевых моделей токи и заряды, сохраняющиеся в силу инвариантности относительно достаточно сложных калибровочных групп, играют ведущую роль.



Выражение (4) для пространственно-временной локализации нетерова тока (это выражение называется ка ноническким) не однозначно, если исходить только из требования выполнения дифференциального закона сохранения (3) и получения правильной интегральной величины (6). Выполнение этих требований не нарушается при замене



\[

J_{A}^{\mu} \rightarrow J_{A}^{\mu}+\frac{\partial f_{A}^{\mu v}\left(\varphi^{a} ; \varphi_{, v}^{a}\right)}{\partial x^{v}} ; \quad f_{A}^{\mu v}+f_{A}^{v \mu}=0

\]



с произвольной функцией $f$. Этим произволом пользуются, чтобы заменить канонич. тензор (5) $T_{\mathrm{v}}^{\mu}$ (не симметричный для отличного от нуля спина) на симметричный (тензор Белинфанте), выбирая



$f^{v, \rho \mu}=\frac{1}{2}\left\{-\frac{\partial L}{\partial \varphi_{, \mu}^{a}} \mathscr{L}^{v \rho ; a b} \varphi^{b}+\frac{\partial L}{\partial \varphi_{, \rho}^{a}} \mathscr{L}^{\mu v ; a b} \varphi^{b}+\frac{\partial L}{\partial \varphi_{, v}^{a}} \mathscr{L}^{\rho \mu ; a b} \varphi^{b}\right\}$.



Для нулевого спина то же преобразование позволяет получить для безмассового поля $T_{\mathrm{v}}^{\mu}$ с нулевым следом.



Однозначные выражения для нетеровых токов получаются варьированием по полям, для к-рых эти токи служат источниками.



Для теорий, обладающих суперсимметрией, независимыми переменными при выводе Н.т. будут наряду с $x$ и антикоммутирующие координаты $\theta_{\alpha}(\alpha-$ спинорный индекс). Это приводит к обобщению фундаментальных сохраняющихся величин, а также к появлению новых сохраняющихся величин: спин-векторных токов и соответствующих им суперзарядов, образующих представление супералгебры Пуанкаре.



Для классич. теорий поля выписанных формальных вы ражений вполне достаточно. В квантовой теории поля выражения (4), (6), как правило, нуждаются в регуляризации и перенормировке. При этом может оказаться, что формально имеющаяся симметрия не может быть сохранена для регуляризованных выражений и соответствуюший закон сохранения перестает выполняться - говорят, что присутствует аномалия. Так, при рассмотрении взаимодействия безмассовых фермионов с электромагнитным полем в классич. теории наряду с векторным сохраняется также и аксиальный ток $j_{5}^{\mu}=\bar{\psi} \gamma^{\mu} \gamma^{5} \psi\left(\gamma^{\mu}, \gamma^{5}-\right.$ матришы Дирака). В кван-



HETEP 385





\end{document}